\begin{corollary}[2.8.2]
\label{IV.2.8.2-fr}
Bajo las condiciones de \textup{(2.8.1)}, para que $\mathcal F$ sea plano respecto de
$f$, es necesario y suficiente que el homomorfismo canónico
$\mathcal F\to i_*(\mathcal F_\eta)=i_*(i^*(\mathcal F))$
sea inyectivo.
\end{corollary}

\begin{env}[2.8.3]
\label{IV.2.8.3-fr}
La formación del $\mathcal O_X$-módulo $\overline{\mathcal G'}$ es
\emph{funtorial en $\mathcal F$ y $\mathcal G'$}: más precisamente, si
$\mathcal F_1$, $\mathcal F_2$ son dos $\mathcal O_X$-módulos
cuasicoherentes, $u:\mathcal F_1\to\mathcal F_2$ un homomorfismo de
$\mathcal O_X$-módulos, $\mathcal G'_i$ un
$\mathcal O_{X_\eta}$-módulo cociente de $(\mathcal F_i)_\eta$ $(i=1,2)$ y
$v:\mathcal G'_1\to\mathcal G'_2$ un homomorfismo que hace conmutativo el
diagrama
\[
\xymatrix{
(\mathcal F_1)_\eta \ar[r]^{i^*(u)} \ar[d] &
(\mathcal F_2)_\eta \ar[d] \\
\mathcal G'_1 \ar[r]_v & \mathcal G'_2
}
\]

\oldpage[IV-2]{35}
(homomorfismo determinado de manera única (cuando existe) por esta
propiedad), entonces el diagrama
\[
\xymatrix{
\mathcal F_1 \ar[r]^u \ar[d] & \mathcal F_2 \ar[d] \\
i_*(\mathcal G'_1) \ar[r]_{i_*(v)} & i_*(\mathcal G'_2)
}
\]
es conmutativo y, por consiguiente, existe un único homomorfismo
$w:\overline{\mathcal G'_1}\to\overline{\mathcal G'_2}$ que hace conmutativo
el diagrama
\[
\xymatrix{
\mathcal F_1 \ar[r]^u \ar[d] & \mathcal F_2 \ar[d] \\
\overline{\mathcal G'_1} \ar[r]_w & \overline{\mathcal G'_2}
}
\]
\end{env}
